\documentclass{article}
\usepackage[utf8]{inputenc}
\usepackage[T1]{fontenc}
\usepackage{geometry}
\usepackage{graphicx}
\usepackage{float}
\usepackage{hyperref}
\usepackage[english,indonesian]{babel}

\title{Optimasi Pemrosesan Template LaTeX Menggunakan Pendekatan Berbasis AI}
\author{Academic Support Agent}
\date{\today}

\begin{document}
\maketitle

\begin{abstract}

\end{abstract}

\chapter{PENDAHULUAN}


\chapter{TINJAUAN PUSTAKA}


% Original content (commented out):
% %% Created by Jonas Thern
%%Version 2.7-Lab
% Weitere Infos in Bedienungshinweise.pdf
\documentclass[a4paper,11pt]{article}
\input{misc/Background/Settings}
\input{misc/Background/Symbolverzeichnis}
\input{Kapitel/00-Symbole}
\input{misc/Background/titlepage}
\input{misc/Setup}
\input{misc/Einheitenmacros}
\addbibresource{Literatur.bib}
\hypersetup{pdfsubject={\Dokumenenart}
             ,pdftitle={\Laborthema}
             ,pdfauthor={\Gruppe}}


%% DOKUMENT 
\begin{document}
\input{misc/Background/Dokumentbeginn}

%% Kapitel ab hier
\input{Kapitel/01-Aufgabe}
\clearpage
\input{Kapitel/02-Theoretische Grundlagen}
\clearpage
\input{Kapitel/03-Versuchsaufbau}
\clearpage
\input{Kapitel/04-Versuchsdurchführung}
\clearpage
\input{Kapitel/05-Versuchsergebnisse}
\clearpage
\input{Kapitel/06-Diskussion}
\clearpage
\printbibliography[heading=bibintoc, title={Literaturverzeichnis}]
\clearpage
\input{Kapitel/100-Anhang}
\end{document}



\end{document}